\section{Introduction}\label{sec:introduction}
Group-level preference facts are expensive to obtain and easy to over-interpret. A proposed survey item must be piloted, translated, and fielded before researchers can assess whether it measures the intended attribute across populations. Language models offer a cheaper preliminary readout: researchers supply an item and a population description and obtain synthetic responses \cite{argyle2023,horton2023,buskirk2026}. Yet sensitivity to prompting, compressed variance, and unstable subgroup comparisons limit the interpretation of those responses \cite{bisbee2024,boelaert2025,wang2025}. A useful methodological question is which population information the researcher explicitly supplies and what subsequent comparisons can establish.

We address this question by constructing \emph{anchored synthetic choice policies}. A researcher declares aggregate preference benchmarks, uses their signs to label a common bank of synthetic contrasts, and fits a language-model adapter with Direct Preference Optimization (DPO) \cite{rafailov2023}. We interpret the loss through a regularized Bradley--Terry pairwise choice model. Each adapter is jointly trained across the preference dimensions. At evaluation, it assigns probabilities to supplied response codes for new items without a country identifier in the prompt. This restriction removes direct country-name conditioning at inference; pretrained knowledge, instruction tuning, generator judgments, and optimization variation remain.

The empirical application uses country scores from the Global Preferences Survey (GPS) \cite{falk2018}. These are aggregates of survey measures selected through prior experimental validation. They are not country-level collections of incentivized experimental choices. We evaluate the fitted policies on candidate indicators in World Values Survey (WVS) Wave~7 \cite{wvs2022}. Three country-level associations distinguish retention of the construction signal from agreement with people: human--anchor, policy--anchor, and policy--human. Held-out synthetic comparisons separately measure whether fitting recovers the imposed labels.

A \emph{construct} is a theoretical preference attribute, a \emph{measure} operationalizes that attribute through answers or choices, and an \emph{anchor} is the resulting aggregate score used in construction. A fitted synthetic policy is a computational scoring instrument. None of these objects is interchangeable. In particular, aggregate scores do not identify a joint distribution of human preferences, and country-level associations do not establish individual-level measurement validity.

Our contributions are threefold:
\begin{enumerate}[leftmargin=1.5em,itemsep=2pt]
\item We specify an inspectable route from declared aggregate preference signs to jointly fitted synthetic choice policies.
\item We evaluate this construction on a separate survey battery, distinguishing synthetic-label recovery, external anchor ordering, and human agreement.
\item We use discrepancies among these readouts to diagnose cross-survey coherence and to identify questions requiring further validation.
\end{enumerate}
The intended application is group-level empirical analysis and survey pre-piloting, as a complement to human research. This application demonstrates text-responsive scoring and trait-dependent anchor transfer. Whether such scores improve item selection before a human pilot remains an empirical question.
